% 勒文海姆–斯科伦定理（综述）
% license CCBYSA3
% type Wiki

本文根据 CC-BY-SA 协议转载翻译自维基百科\href{https://en.wikipedia.org/wiki/L\%C3\%B6wenheim\%E2\%80\%93Skolem_theorem}{相关文章}。

在数理逻辑中，洛文海姆–斯科伦定理（Löwenheim–Skolem theorem）是关于模型的存在性与基数的一个定理，以数学家莱奥波德·洛文海姆（Leopold Löwenheim）和托拉夫·斯科伦（Thoralf Skolem）的名字命名。  

其精确定义如下：若一个可数的一阶理论（countable first-order theory）存在一个无限模型，那么对于任意无限基数 \(\kappa\)，该理论都存在一个基数为 \(\kappa\) 的模型。这意味着，任何具有无限模型的一阶理论都不可能在同构意义下拥有唯一模型。  

由此可知，一阶理论无法控制其无限模型的基数。  

（向下的）洛文海姆–斯科伦定理是刻画一阶逻辑的两个关键性质之一，另一个是紧致性定理**（compactness theorem）。这两者共同构成了**林德斯特伦定理（Lindström’s theorem）的基础。而在更强的逻辑体系中（例如**二阶逻辑**），一般而言洛文海姆–斯科伦定理并不成立。[1]  
```
